\documentclass{template/socthesis}
\addbibresource{text.bib}

\usepackage{subcaption}
\usepackage{amsmath}

\setcounter{secnumdepth}{4}
\setcounter{tocdepth}{4}

\titlecz{Coopmaster - systém pro řízení a automatizaci kurníku}
\titleen{Coopmaster - system for coop control and automation}
\author{Jaroslav Němec}
\field{18}
\school{Střední škola a vyšší odborná škola aplikované kybernetiky}
\mentor{Ing. David Podzimek}
\mentorstatement{Ing. Davida Podzimka}
\newcommand{\mentorthanks}{Ing. Davidu Podzimkovi}

\begin{document}

    \maketitle

    \makecopyrightstatement{V~Hradci Králové}

    \makethanks{Děkuji svému školiteli \mentorthanks{} za obětavou pomoc, podnětné připomínky a nekonečnou trpělivost, kterou mi během práce poskytoval.}

    \pagestyle{empty}

    \section*{Anotace}
    Ve své práci jsem se zabýval návrhem a realizací kompletního řešení pro malé i střední zemědělce a hospodáře, které má pomoci plnit běžné činnosti.
    Cílem této práce bylo vytvořit a realizovat systém pro kontrolu a automatizaci chovu hospodářských zvířat.

    \subsection*{Klíčová slova}
    programování, automatizace, zemědělství, neuronové sítě, python, mikroservisní architektura

    \vspace{20mm}

    \section*{Annotation}
    In my work, I have been involved in the design and implementation of a complete solution for small and medium farmers and homesteaders to help them carry out their day-to-day activities.
    The aim of this work was to design and implement a system for control and automation of livestock farming.

    \subsection*{Keywords}
    programming, automation, agriculture, neural networks, python, microservice architecture

    \pagestyle{plain}

    \tableofcontents % vysází obsah

    %%% Začátek práce
    \setcounter{figure}{0}
    \setcounter{table}{0}
%    \newpage

    %%% Úvod
    \chapter{Úvod}
V dnešní době je populární využívat moderní technologie a umělou inteligenci v různých aplikacích jako například monitoring pohybu zákazníků v obchodech nebo ve formě textových modelů čímž je například známe ChatGPT. Zároveň dnes ubývá lidí, kteří by se chtěli zabývat staráním se o hospodářská zvýřata a tudíž je potřeba, aby něco zaujmulo jejich místo. Zároveň jsem sám člověk, který rád zkoumá nové věci v oblasti informatiky a velice ho baví automatizování nejrůznějších úloh. Na základě těchto faktů jsem se rozhodl vytvořit téma práci popisující využití moderních technologií při chovu hospodářských zvýřat. Nápad na tuto práci vznikl díky mojí babičce, která chová doma slepice. Když odjede na dovolenou, stávám se já tím, kdo se o ně musí starat. Slepice je třeba dojet ráno a večer zkontrolovat a spočítat. Tato činnost je časově náročná kvůli cestování a kolikrát i zbytečná, protože většinou se nic neděje. A vzhledem k tomu, že jsem povoláním informatik, budu toto práci věnovat tomu jaké je moje řešení automatizace babiččina chovu.
\newline
Práce se dělí na teoretickou a praktickou část. V teoretické části jsou popsány základní pojmy a principy, jež jsou důležité pro plné chápání práce. Praktická část popisuje konkrétní implementaci a nasazení asistenčního systému pro chov hospodářských zvýřat. Čtenář může využít znalosti, které nasbírá během čtení teoretické části a jež mu následně pomůže chápat praktická část, jako návod k realizaci vlastního systému dle jeho potřeb.
\newline
Svojí konkrétní implementaci jsem zaměřil na chov Kura domácího z výše popsaných osobních důvodů, a protože byl pro mě nejdostupnější testovací zvíře. Systém jsem realizoval mikroservisní architekturou a jednotlivé služby jsou psány v populárnímu jazyce Python. Jako GUI rozhraní je šikovně použit open source software pro chytrou domácnost Home Assistant.
\newline
Ukázkový běh aplikace je dostupný na url adrese https://coopmaster.jarousnemec.cz/ s přihlašovacími údaji do Home Assistanta uživatel: \textbf{visitor} a heslo: \textbf{Heslo1234}. Zdrojový kód k jednotlivým službám a konfiguraci Home Assistanta je dostupný na serveru github.com.


    \chapter{Teoretická část}
\section{Kontejnerizace a docker compose}
\section{Mikroservisní architektura}
\section{Jazyk Python}
\section{Jazyky Wire a C++ }
\section{Git}
\section{Github workflow}
\section{Mqtt}
\section{Home Assistant}
\section{Arduino}
\section{Ip kamera a RTSP}
\section{Cloudflare tunneling}
\section{Flask}
\section{Strojové učení}
\section{Yolo Ultralytics}
\section{Tenzometrický senzor}
\section{Poe}
\section{Wifi extender}
\section{TailScale vpn}
\section{Raspberry PI}
\section{Raspberry PI OS}
\section{Yaml}
\section{GUI}






    \chapter{Praktická část}\label{ch:prakticka2-cast}
\section{Podvozek}\label{sec:podvozek}
\subsection{Problémy}\label{subsec:podvozek-problemy}
\begin{itemize}
    \item potřebovali sme internetové připojení do kurníku
    \item elektriku
    \item vpn
    \item build docker imagu
    \item gui
    \item silný výpočetní výkon
\end{itemize}


\subsection{Možnosti řešení}\label{subsec:podvozek-moznosti-reseni}
\subsection{Řešení problémů}\label{subsec:podvozek-reseni-problemu}
\subsection{Zhodnocení}\label{subsec:podvozek-zhodnoceni}
\section{Chytrá dvířka}\label{sec:chytra-dvirka}
\subsection{Stav před}\label{subsec:chytra-dvirka-problemy}
-babi jezdí na dovolenou musím jezdit já\newline
-ráno zapomenu nebo nepomene pustit\newline
-večer zapomenem zavřít\newline
-děda nepočítá slepice\newline
\subsection{Možnosti řešení}\label{subsec:chytra-dvirka-moznosti-reseni}
Automatizace zavírání kurníku představuje praktické řešení pro chovatele drůbeže, kteří chtějí zajistit bezpečí svých zvířat před predátory a zároveň si ulehčit každodenní rutinu. Existuje několik přístupů k realizaci tohoto úkolu, z nichž zde uvedeme tři potenciální řešení: zakoupení hotového produktu, vývoj vlastního řešení s využitím kladky a táhlového motoru a vývoj s využitím krokového motoru.\newline

1. Zakoupení hotového řešení\newline
Prvním a nejjednodušším řešením je zakoupení již existujícího automatizovaného systému na internetu. Tento přístup nabízí několik výhod:

Snadná instalace: Komerčně dostupné produkty jsou navrženy s ohledem na snadnou montáž. Uživatelé obvykle obdrží podrobný manuál, který jim pomůže s instalací systému.
Časová úspora: Zakoupení hotového produktu šetří čas, který by jinak byl potřeba na vývoj a testování vlastního řešení.
Záruka a podpora: Výrobci často poskytují záruku a technickou podporu, což zvyšuje jistotu spolehlivosti systému.
Na druhou stranu, nevýhodou může být omezená možnosti přizpůsobení specifickým potřebám uživatele nebo vyšší pořizovací náklady.\newline

2. Vývoj vlastního řešení s kladkou a táhlovým motorem\newline
Druhé řešení zahrnuje vývoj vlastního systému využívajícího kladky a táhlový motor. Tento přístup umožňuje větší flexibilitu a přizpůsobení se konkrétním podmínkám kurníku:

Přizpůsobitelnost: Navržený systém může být přizpůsoben velikosti a tvaru dveří kurníku.
Nízké náklady: V závislosti na dostupnosti komponentů může být tento přístup ekonomicky výhodnější než komerční řešení.
Praktické dovednosti: Uživatel získá technické dovednosti spojené s vývojem a implementací mechanických systémů.
Nevýhodou tohoto přístupu je časová náročnost spojená s návrhem a testováním systému, a potenciální potřeba seznámit se s technickými detaily týkajícími se motorů a kladek.\newline

3. Vývoj vlastního řešení s krokovým motorem\newline
Třetí možností je návrh vlastního systému využívajícího krokový motor. Tento přístup zahrnuje pokročilější technické znalosti, ale nabízí možnosti přesné kontrolování pohybu dveří kurníku:

Přesnost: Krokový motor umožňuje přesnou kontrolu rychlosti a polohy dveří, což je užitečné v případě potřeby plynulého zavírání a otevírání.
Programovatelnost: Systém může být naprogramován k automatickému otevírání a zavírání v závislosti na čase, světelných podmínkách nebo jiných parametrech.
Možnosti rozšíření: Toto řešení může být rozšířeno o další senzory a kontrolní mechanismy pro zvýšení funkčnosti.
Na druhou stranu, tento přístup vyžaduje více času na vývoj a testování, a také znalosti v oblasti elektroniky a programování.

Každé z těchto řešení má své výhody a nevýhody, a výběr závisí na specifických potřebách uživatele, jeho technických schopnostech a rozpočtu.


slepice zůstanou venku
Bezpečné zavírání dvířek kurníku je klíčové, aby nedocházelo k uzavření slepic mimo kurník, kde by byly ohroženy predátory nebo nepříznivými povětrnostními podmínkami. Zde je několik tipů a technik, jak toho dosáhnout:

Použití senzorů pohybu nebo přítomnosti:

Instalujte senzory pohybu nebo infrapasivní senzory (PIR), které detekují přítomnost slepic v blízkosti dvířek. Pokud jsou senzory aktivovány, dvířka zůstanou otevřená nebo se znovu otevřou, aby umožnily slepicím bezpečný vstup.
Časové zpoždění zavírání:

Nastavte systém na časové zavírání dvířek s dostatečným zpožděním po soumraku, aby se všechny slepice stihly vrátit do kurníku. Pomocí časovače můžete určit optimální čas zavírání na základě ročních období.
Světelný senzor:

Použití světelného senzoru, který reaguje na změnu světla při západu slunce, umožní dvířkům se automaticky zavřít. Tento systém může být kalibrován tak, aby zohlednil přirozené chování slepic vracících se do kurníku, když se stmívá.
Ruční ověření:

Pokud je to možné, provádějte před zavřením dvířek poslední vizuální kontrolu, abyste se ujistili, že jsou všechny slepice v kurníku.
Zvukové signály:

Před zavřením dvířek může být spuštěn zvukový signál, jako je bzučák nebo jemný alarm, který slepicím naznačí, že se dvířka brzy zavřou, což je přiměje vstoupit do kurníku.
Postupné zavírání:

Počítání pomocí kamery

Implementujte systém, který dveře zavírá postupně, což minimalizuje riziko, že by slepice zůstaly uvězněné venku. Tento systém může fungovat na principu pomalého sklápění dveří pomocí krokového motoru.
Použití těchto opatření může výrazně snížit riziko, že by slepice zůstaly venku, a zároveň zvýšit celkovou bezpečnost a efektivitu provozu automatizovaného systému zavírání dvířek kurníku.

\subsection{Řešení}\label{subsec:chytra-dvirka-reseni-problemu}
Vytvoření automatických dvířek řízených Arduinem s pohonem na táhlový motor a integrací s Home Assistantem prostřednictvím MQTT může být zajímavý projekt. Zde je podrobný postup, jak můžete tento projekt uskutečnit:

Materiály a nástroje
Konstrukční materiál:

Hliníkové profily a úhelníky pro rám a dvířka
Spojovací materiál (šrouby, matice)
Elektronika:

Arduino (např. Arduino Uno)
H-můstek (např. L298N) pro ovládání motoru
Táhlový motor (např. stejnosměrný motor s posuvníkem)
Napájecí adaptér pro motor (podle specifikace motoru)
Senzory a další komponenty:

Koncové spínače pro detekci polohy dvířek (otevřené/zavřené)
Modul ESP8266 nebo ESP32 pro Wi-Fi připojení a komunikaci s MQTT
Breadboard a propojovací vodiče
Software:

Arduino IDE
MQTT broker (např. Mosquitto)
Home Assistant upravený pro integraci MQTT
Mobilní aplikace pro Home Assistant
Krok za krokem
Mechanická konstrukce:
Návrh a sestavení:

Navrhněte rozměry dvířek a rámu s ohledem na montáž táhlového motoru a připojení koncových spínačů.
Sestavte hliníkový rám a upevněte pohyblivou část ke konstrukci pomocí táhlového motoru.
Montáž táhlového motoru:

Upevnite táhlový motor tak, aby mohl otevírat a zavírat dvířka.
Připevněte koncové spínače, které budou detekovat koncové polohy (úplně otevřeno, úplně zavřeno).
Elektronická konstrukce:
Zapojení elektroniky:
Připojte motor k H-můstku a ten k Arduinu, společně s napájením.
Zapojte koncové spínače k Arduinu pro detekci polohy dvířek.
Přidejte ESP8266/ESP32 k Arduinu pro Wi-Fi komunikaci.
Naprogramování Arduina:
Arduino kód:
Naprogramujte Arduino k přijímání příkazů přes sériovou linku z ESP8266/ESP32 (pomocí knihovny pro komunikaci s MQTT).
Implementujte logiku pro řízení motoru na základě příkazů (otevřít/zavřít), s využitím údajů z koncových spínačů pro bezpečné zastavení motoru.
Použijte H-můstek pro přepínání polarity motoru, což umožní reverzaci směru otáčení.
Integrace s MQTT:
Nastavení MQTT:

Zprovozněte MQTT broker. Pokud používáte Mosquitto, můžete ho spustit na serveru nebo Raspberry Pi.
Nastavte ESP8266/ESP32 pro odesílání a přijímání MQTT zpráv z brokera.
Integrace s Home Assistant:

Přidejte MQTT integraci do Home Assistantu a vytvořte entity pro řízení dvířek (pomocí YAML konfigurace).
Nastavte automatizace nebo scény, pokud je chcete ovládat na základě dalších podmínek (např. časovač, senzor).
Ovládání z mobilní aplikace:
Nastavení mobilní aplikace:
Připojte mobilní aplikaci Home Assistant k vaší instanci Home Assistantu.
Vytvořte uživatelsky přívětivé rozhraní pro ovládání dvířek z aplikace.
Testování a ladění:
Testování funkčnosti:
Ověřte správnost mechanické instalace a zkontrolujte, zda dvířka správně reagují na příkazy.
Testujte komunikaci přes MQTT a přesvědčte se, že příkazy z mobilu správně aktivují činnost dvířek.
Otestujte automatizace a jejich spolehlivost.
Tento projekt kombinuje mechanickou konstrukci, elektroniku a softwarové řešení. Ujistěte se, že každou část projektu pečlivě testujete a ladíte pro zajištění bezpečného a spolehlivého provozu.


jak jsem to konkrétně realizoval

ovladaní dvířek
bezpečné zavírání - počítaní
\subsection{Zhodnocení}\label{subsec:chytra-dvirka-zhodnoceni}

Při zhodnocení nového stavu automatických dvířek postavených podle navrženého řešení je důležité zvážit obě strany – výhody i nevýhody.

Klady
Automatizace a pohodlí:

Dvířka se otevírají a zavírají automaticky, což zvyšuje pohodlí, zejména pokud máte často obsazené ruce.
Dálkové ovládání:

Možnost ovládat dvířka na dálku pomocí mobilní aplikace představuje pohodlné a flexibilní řešení pro uživatele.
Integrace s Home Assistantem:

Systém je kompatibilní s Home Assistantem, což umožňuje integraci do širšího ekosystému chytré domácnosti a automatizace různých úloh.
Možnost programování:

Díky použití Arduina je možné systém snadno programovat a přizpůsobovat funkcionalitu podle potřeby.
Vzdálené monitorování a ovládání:

Schopnost monitorovat stav dvířek a řídit jejich funkčnost přes MQTT zvyšuje míru kontroly a bezpečnosti.
Zápory
Komplexita instalace:

Sestavení a instalace mohou být složité a časově náročné, vyžadují znalosti z elektroniky, programování a konstrukce.
Náklady:

I když jsou komponenty jako Arduino a ESP8266 relativně levné, celkové náklady mohou narůst včetně motoru a dalších komponentů.
Spolehlivost:

Systém může být náchylný k problémům spojeným s výpadky připojení, chybami v softwaru nebo mechanickými poruchami.
Spotřeba energie:

Systém je závislý na elektrické energii, takže v případě výpadku proudu může dojít k nefunkčnosti.
Bezpečnost a zabezpečení:

Nezabezpečený přístup k MQTT brokeru nebo k systému Home Assistant by mohl způsobit bezpečnostní rizika.
Údržba:

Systém může vyžadovat pravidelnou údržbu, zejména u mechanických částí, aby bylo zajištěno jejich správné fungování.
Závěr
Realizace tohoto projektu může významně zvýšit komfort a automatizaci ve vaší domácnosti, ale přinese s sebou i určité výzvy, které je potřeba zvážit a připravit se na jejich řešení. Úspěch projektu bude záviset na kvalitní realizaci, testování a následné údržbě.
\section{Detekce snesených vajec}\label{sec:detekce-snesenych-vajec}
\subsection{Problémy}\label{subsec:detekce-snesenych-vajec-problemy}
\subsection{Možnosti řešení}\label{subsec:detekce-snesenych-vajec-moznosti-reseni}
\subsection{Řešení problémů}\label{subsec:detekce-snesenych-vajec-reseni-problemu}
\subsection{Zhodnocení}\label{subsec:detekce-snesenych-vajec-zhodnoceni}
\section{Detekce vetřelců}\label{sec:detekce-vetrelcu}
\subsection{Problémy}\label{subsec:detekce-vetrelcu-problemy}
\subsection{Možnosti řešení}\label{subsec:detekce-vetrelcu-moznosti-reseni}
\subsection{Řešení problémů}\label{subsec:detekce-vetrelcu-reseni-problemu}
\subsection{Zhodnocení}\label{subsec:detekce-vetrelcu-zhodnoceni}
\section{Chytré světlo}\label{sec:chytre-svetlo}
\subsection{Problémy}\label{subsec:chytre-svetlo-problemy}
\subsection{Možnosti řešení}\label{subsec:chytre-svetlo-moznosti-reseni}
\subsection{Řešení problémů}\label{subsec:chytre-svetlo-reseni-problemu}
\subsection{Zhodnocení}\label{subsec:chytre-svetlo-zhodnoceni}
\section{Monitoring podmínek}\label{sec:monitoring-podminek}
\subsection{Problémy}\label{subsec:monitoring-podminek-problemy}
%nechat teplotu a vlhkost
\subsubsection{Ovládatní dvířek, snímání teploty a vlhkosti}
Pro ovládatní dvířek a snímání povětrnostních podmínek je nejrozumnější vytvořit řídící jednotku, která tyto požadavky obsáhne.
Další částí částí jsou samotná elektrická dvířka, jejichž funkce bude na základě způsobu elektrického signálu být otevřená nebo zavřená.
Teplota a vlhkost jsou snímány senzory, jež jsou souřástí řídící jednotky.
Implementace je popsána v sekci~\ref{subsec:ridici-jednotka}.
\subsection{Možnosti řešení}\label{subsec:monitoring-podminek-moznosti-reseni}
\subsection{Řešení problémů}\label{subsec:monitoring-podminek-reseni-problemu}
\subsection{Zhodnocení}\label{subsec:monitoring-podminek-zhodnoceni}





    \chapter{Zamyšlení nad možnými dalšími příležitostmi pro automatizaci}

\section{Co vše se dá automatizovat?}
\begin{itemize}
    \item skleník
    \item akvarium
    \item kravín a pastviny
    \item záhony a zahrada
\end{itemize}

\section{Automatizace skleníku}
- kontrola vlhkosti, větrání na základě teploty ve skleníku, přidávání hnojiva do zálivky

\section{Chytré akvárium}
- chytré topení, filtrace, krmítko, kontrola kvality vody

\section{Monitoring kravína a pastvin}
- zranění kráva\newline
- cizí pes\newline
- cizí člověk

\section{Monitoring záhonů}
- pes krade mrkev\newline
- utekly slepice\newline
- zajíc\newline
- plíseň na bramborách\newline
- sucho

    %%% Závěr
    \newpage
\chapter{Závěr}
Závěrem lze říct, že jsme úspěšně vyvinuli systém, který je schopný v reálném čase pomáha plnit každodení úkony, které by hospodář musel jinak dělat sám. Díky naší aplikaci může člověk kontrolovat situaci v kurníku případně i ve výběhu. Systém ho sám informuje, pokud je něco v nepořádku, například v případě, že je ve výběhu vetřelec, který by mohl představovat nebezpečí pro chované slepice. Dále se nám na podobném principu povedlo implementovat kontrolu počtu slepic v kurníku, což lze použít například při automatizaci zavírání dvířek. Systém také poskytuje aktuální záběry z bezpečnostních kamer díky, kterým si může hospodář sám kontrolovat situaci svého chovu a nemusí díky tomu být přítomen fyzicky. V poslední řadě jsme také implementovali část systému, která je schopna kontrolovat stavy jednotlivých hnízd a zjistit zda v něm sedí slepice případně kolik vajec v hnízdě je. Tato data lze také použít ke tvorbě statistik, díky nimž může člověk analizovat různé aspekty svého chovu.

\begin{figure}[h]
    \centering
    \includegraphics[width=\textwidth]{img/pulp}
    \caption{Pulp Fiction}
    \label{fig:pulp-fiction}
\end{figure}



    \printbibliography[title=Použitá Literatura]
    \addcontentsline{toc}{section}{Použitá Literatura}

    \listoffigures
    \addcontentsline{toc}{section}{Seznam obrázků}

    \listoftables
    \addcontentsline{toc}{section}{Seznam tabulek}

    \listoflistedequation
    \addcontentsline{toc}{section}{Seznam rovnic}

\end{document}
